\section*{Conclusions}

The threshold-defined primary analysis comprised 68 analyzable tests. TYK2 was the only FDR-significant gene and served as an established psoriasis-locus positive control. Colocalization supported a shared signal in sun-exposed skin and whole blood but not in non-sun-exposed skin. Conditional cis-eQTL coverage within the GTEx significant-pair files was limited to 18 of 57 gene--tissue cells, and only 17 cells could be analyzed against all four outcomes. A relaxed-threshold PDE4A association lacked locus-level colocalization support and remains exploratory. These findings emphasize that instrument scarcity, single-variant uncertainty, and the distinction between transcript effects and pharmacologic modulation must be explicit in dermatologic drug-target MR.

\section*{Abbreviations}

cis-eQTL, cis-expression quantitative trait locus; CI, confidence interval; FDR, false discovery rate; GTEx, Genotype-Tissue Expression; GWAS, genome-wide association study; LD, linkage disequilibrium; MR, Mendelian randomization; OR, odds ratio; PP.H2, posterior probability of association with the GWAS trait only; PP.H3, posterior probability of association with both traits through distinct causal variants; PP.H4, posterior probability of a shared causal variant; SE, standard error; SNP, single nucleotide polymorphism; STROBE-MR, Strengthening the Reporting of Observational Studies in Epidemiology using Mendelian Randomization; TYK2, tyrosine kinase 2.
